% ============================================================
%  EXAMEN TEMA 1: ARITMÉTICA CON NÚMEROS NATURALES Y DECIMALES
%  Curso Introductorio de Matemáticas — 1er Año
%  Autor: Ing. Gabriel Astudillo
%  Sesiones cubiertas: Semana 1 (clases 1 a 4)
%  Nota: enunciado limpio (sin pistas ni desarrollo). Las
%  soluciones están en la "Clave de Respuestas" al final.
% ============================================================

\documentclass[12pt, letterpaper]{article}

% ── Paquetes ─────────────────────────────────────────────────

\usepackage{mathwizards-guia}
\mwguia{Examen — Tema 1: Aritmética con Decimales}{Ing. Gabriel Astudillo $\cdot$ Curso 1er Año}

% ── Comandos propios de este examen ───────────────────
\newcommand{\Z}{\mathbb{Z}}

\begin{document}
% ============================================================

% ── Portada / Título ─────────────────────────────────────────
\mwportada{Examen del Tema 1}{Aritmética con Números Naturales y Decimales}{Sumar, Restar, Multiplicar y Dividir con Rigor Posicional}

\vspace{4pt}
\begin{cuadroinfo}
  \small
  \textbf{Nombre:} \rule{0.55\linewidth}{0.4pt} \quad \textbf{Fecha:} \rule{0.15\linewidth}{0.4pt} \\[6pt]
  \textbf{Tiempo sugerido:} 60 minutos \quad$\bullet$\quad \textbf{Puntaje total:} 100 puntos \\[4pt]
  \textbf{Instrucciones:}
  \begin{enumerate}[leftmargin=*]
    \item Lee cada ejercicio con calma antes de resolverlo.
    \item Muestra \emph{todos} tus pasos: se evalúa cómo llegas a la respuesta, no solo la respuesta.
    \item Presta mucha atención a la \emph{posición} de la coma decimal y a los \emph{ceros} no significativos.
    \item No se permite el uso de calculadora.
    \item Escribe con letra clara y revisa tu examen antes de entregarlo.
  \end{enumerate}
\end{cuadroinfo}

\vspace{8pt}

% ============================================================
\section{Suma con Decimales}
% ============================================================

\begin{enumerate}[leftmargin=*]
  \item Calcule el valor exacto:
        \[
          99{,}9999 + 0{,}00109
        \]

  \item Obtenga la suma sin redondear:
        \[
          301{,}005 + 0{,}03105
        \]

  \item Determine el valor exacto:
        \[
          919{,}0909 + 0{,}90809
        \]

  \item Efectúe la siguiente suma:
        \[
          1989{,}0099 + 0{,}890199
        \]
\end{enumerate}

\newpage

% ============================================================
\section{Resta con Decimales}
% ============================================================

\begin{enumerate}[leftmargin=*]
  \item Resuelva la sustracción:
        \[
          10000{,}00001 - 998{,}99899
        \]

  \item Realice la siguiente operación:
        \[
          305{,}1 - 148{,}9980706
        \]

  \item Calcule:
        \[
          2000{,}003 - 1987{,}007985
        \]

  \item Resuelva la sustracción (préstamo sobre el cero):
        \[
          1 - 0{,}999
        \]
\end{enumerate}

\newpage

% ============================================================
\section{Multiplicación con Decimales}
% ============================================================

\begin{enumerate}[leftmargin=*]
  \item Calcule el producto exacto:
        \[
          0{,}025 \times 0{,}0040
        \]

  \item Determine el valor exacto:
        \[
          1{,}0238 \times 0{,}00075
        \]

  \item Efectúe la multiplicación exacta:
        \[
          99{,}0013 \times 0{,}0099
        \]

  \item Calcule el producto (ceros devorados):
        \[
          0{,}1254 \times 0{,}008
        \]
\end{enumerate}

\newpage

% ============================================================
\section{La Unidad Seguida de Ceros}
% ============================================================

\begin{enumerate}[leftmargin=*]
  \item Calcule (mueva la coma a la derecha):
        \[
          0{,}0035 \times 10000
        \]

  \item Calcule (mueva la coma a la derecha):
        \[
          0{,}00789 \times 10000
        \]

  \item Calcule (mueva la coma a la izquierda):
        \[
          50 \div 10000
        \]

  \item Calcule (mueva la coma a la izquierda):
        \[
          78973{,}5 \div 100
        \]
\end{enumerate}

\newpage

% ============================================================
\section{División con Decimales}
% ============================================================

\begin{enumerate}[leftmargin=*]
  \item Calcule (decimal entre decimal):
        \[
          0{,}75 \div 0{,}05
        \]

  \item Calcule (entero entre decimal):
        \[
          450 \div 0{,}6
        \]

  \item Calcule (decimal entre decimal):
        \[
          15\,625 \div 0{,}125
        \]

  \item Calcule hasta dos cifras decimales (decimal entre entero):
        \[
          25{,}5 \div 7
        \]
\end{enumerate}

\newpage

% ============================================================
\ifmwclaves
  \section{Clave de Respuestas}
  % ============================================================

  \subsection*{Suma con Decimales (Ejercicios 1--4)}

  \begin{enumerate}[leftmargin=*, label=\textbf{\arabic*.}]
    \item $99{,}9999 + 0{,}00109 = 100{,}00099$
    \item $301{,}005 + 0{,}03105 = 301{,}03605$
    \item $919{,}0909 + 0{,}90809 = 919{,}99899$
    \item $1989{,}0099 + 0{,}890199 = 1989{,}900099$
  \end{enumerate}

  \subsection*{Resta con Decimales (Ejercicios 1--4)}

  \begin{enumerate}[leftmargin=*, label=\textbf{\arabic*.}]
    \item $10000{,}00001 - 998{,}99899 = 9001{,}00102$
    \item $305{,}1 - 148{,}9980706 = 156{,}1019294$
    \item $2000{,}003 - 1987{,}007985 = 12{,}995015$
    \item $1 - 0{,}999 = 0{,}001$
  \end{enumerate}

  \subsection*{Multiplicación con Decimales (Ejercicios 1--4)}

  \begin{enumerate}[leftmargin=*, label=\textbf{\arabic*.}]
    \item $0{,}025 \times 0{,}0040 = 0{,}0001$ \quad ($10^{-4}$)
    \item $1{,}0238 \times 0{,}00075 = 0{,}00076785$
    \item $99{,}0013 \times 0{,}0099 = 0{,}98011287$
    \item $0{,}1254 \times 0{,}008 = 0{,}0010032$ \quad (pues $1254 \times 8 = 10032$, $7$ decimales)
  \end{enumerate}

  \subsection*{La Unidad Seguida de Ceros (Ejercicios 1--4)}

  \begin{enumerate}[leftmargin=*, label=\textbf{\arabic*.}]
    \item $0{,}0035 \times 10000 = 35$
    \item $0{,}00789 \times 10000 = 78{,}9$
    \item $50 \div 10000 = 0{,}005$
    \item $78\,973{,}5 \div 100 = 789{,}735$
  \end{enumerate}

  \subsection*{División con Decimales (Ejercicios 1--4)}

  \begin{enumerate}[leftmargin=*, label=\textbf{\arabic*.}]
    \item $0{,}75 \div 0{,}05 = 75 \div 5 = 15$
    \item $450 \div 0{,}6 = 4500 \div 6 = 750$
    \item $15\,625 \div 0{,}125 = 15\,625\,000 \div 125 = 125\,000$
    \item $25{,}5 \div 7 = 3{,}64$ \quad (resto $0{,}02$)
  \end{enumerate}

  \vspace{10pt}
  \begin{cuadrosolucion}
    \textbf{\textquestiondown Cómo te fue?} Revisa especialmente la \emph{posición} de la coma y los \emph{ceros no significativos}: son la clave para dominar este tema. Si algo se te complicó, vuelve a la Guía 1 y practica el rigor posicional.
  \end{cuadrosolucion}

\fi
\end{document}
